\section{Minkowski Structure and the Soldering Map}
\label{sec:app-lorentz:main}
\label{sec:app-soldering:main}

This appendix establishes, with explicit mathematics and consistent epistemic labels, the chain
by which a Lorentzian four-manifold is \emph{soldered} onto the
$E_{6}$/Albert-algebra arena. We prove the algebraic facts
$H_{2}(\mathbb{O})\cong\mathbb{R}^{1,9}$ and $SL(2,\mathbb{O})\cong\mathrm{Spin}(1,9)$
as theorems; we state the choice of a four-plane, the dimension $d=4$, and the
signature $(1,3)$ as the soldering postulate (P3); and we give the
Osterwalder--Schrader/Page--Wootters route to Lorentzian unitary dynamics in place
of any ``$e^{0}\mapsto i\,e^{0}$'' Wick trick.

The governing logical separation, enforced throughout, is the following. The
\emph{positive-definite} object is the coset/target metric $K|_{\mathfrak{m}}$ on
the $32$-dimensional internal tangent space
$\mathfrak{m}=\mathbf{16}_{-3}\oplus\overline{\mathbf{16}}_{+3}$ of
$M=E_{6}/(\mathrm{Spin}(10)\times U(1))=\mathrm{EIII}$ (see
\autoref{cor:VacuumManifold} and \autoref{thm:uniqueMetric}). The \emph{indefinite}
object is the cubic norm $N$ on the matter representation $\mathbf{27}=J_{3}(\mathbb{O})$,
whose $2\times2$ octonion-Hermitian sub-block carries a Minkowski form on the
\emph{different} space $\mathbb{R}^{1,9}$. No single form is asked to be both
definite and indefinite; the soldering map (P3) is the bridge importing the
indefinite $(1,3)$ structure onto a four-plane of the world-tangent bundle, leaving
$K|_{\mathfrak{m}}$ untouched and Euclidean. Conventions are fixed globally:
signature $(-,+,+,+)$; the dynamics live on the \emph{compact} real form of $E_{6}$
(maximal compact $H=\mathrm{Spin}(10)\times U(1)$), and $E_{6(-26)}$ enters only as
the rigid orbit-classifier of the real $\mathbf{27}$.

\subsection{The Minkowski block: \texorpdfstring{$H_{2}(\mathbb{O})\cong\mathbb{R}^{1,9}$}{H2(O) = R(1,9)}}
\label{subsec:app-lorentz:h2o}

Let $\mathbb{O}$ be the real division octonions, a non-associative composition
algebra of real dimension $8$, with conjugation $x\mapsto\bar{x}$, norm
$n(x)=x\bar{x}=\bar{x}x\in\mathbb{R}_{\ge0}$, and trace $t(x)=x+\bar{x}\in\mathbb{R}$.
Define the space of $2\times2$ octonion-Hermitian matrices
\begin{equation}
\label{eq:app-lorentz:h2def}
H_{2}(\mathbb{O})
 =\left\{\,X=\begin{pmatrix} a & x \\ \bar{x} & b \end{pmatrix}
   \;:\; a,b\in\mathbb{R},\ x\in\mathbb{O}\,\right\},
 \qquad X^{\dagger}=X,
\end{equation}
so that as a real vector space
$\dim_{\mathbb{R}}H_{2}(\mathbb{O})=1+1+8=10$. The determinant is well defined
despite non-associativity (the off-diagonal entry and its conjugate associate with
the scalar diagonal):
\begin{equation}
\label{eq:app-lorentz:det}
\det X = ab-x\bar{x}=ab-n(x).
\end{equation}

\begin{theorem}[Minkowski identification; \texorpdfstring{$H_{2}(\mathbb{O})\cong\mathbb{R}^{1,9}$}{H2(O)=R(1,9)}]
\label{thm:app-lorentz:Mink}
Introduce light-cone-adapted coordinates $a=X^{0}+X^{9}$, $b=X^{0}-X^{9}$, and write
$x=(X^{1},\dots,X^{8})\in\mathbb{R}^{8}$ in an orthonormal basis of $\mathbb{O}$, so
that $n(x)=\sum_{i=1}^{8}(X^{i})^{2}$. Then
\begin{equation}
\label{eq:app-lorentz:detmink}
\det X = (X^{0})^{2}-(X^{9})^{2}-\sum_{i=1}^{8}(X^{i})^{2}
       = -\,\eta_{\mu\nu}\,X^{\mu}X^{\nu},
\qquad
\eta=\mathrm{diag}(-1,\underbrace{+1,\dots,+1}_{9}),
\end{equation}
so the quadratic form $-\det$ on the ten-dimensional space $H_{2}(\mathbb{O})$ has
Lorentzian signature $(1,9)$ in the $(-,+,\dots,+)$ convention:
$(H_{2}(\mathbb{O}),-\det)\cong\mathbb{R}^{1,9}$. \emph{[Theorem.]}
\end{theorem}

\begin{proof}
By \eqref{eq:app-lorentz:det},
$\det X=ab-n(x)=(X^{0}+X^{9})(X^{0}-X^{9})-\sum_{i}(X^{i})^{2}
=(X^{0})^{2}-(X^{9})^{2}-\sum_{i}(X^{i})^{2}$, which is the displayed Minkowski
form. By Sylvester's law of inertia the diagonal pair contributes one $+$ (the
$X^{0}$ mode) and one $-$ (the $X^{9}$ mode), and $-n(x)$ contributes eight further
$-$'s; hence $\det$ has signature $(1,9)$ and $-\det$ is the single-timelike
$\mathbb{R}^{1,9}$.
\end{proof}

\begin{theorem}[Octonionic exceptional isomorphism; \texorpdfstring{$SL(2,\mathbb{O})\cong\mathrm{Spin}(1,9)$}{SL(2,O)=Spin(1,9)}]
\label{thm:app-lorentz:SL2O}
There is a Lie-group isomorphism $SL(2,\mathbb{O})\cong\mathrm{Spin}(1,9)$, under
which the defining $10$-dimensional vector representation acts on
$(H_{2}(\mathbb{O}),-\det)\cong\mathbb{R}^{1,9}$ and the two chiral
$16$-dimensional spinor representations are realised on
$\mathbb{O}^{2}\cong\mathbb{R}^{16}$. This is the octonionic ($n=8$) member of the
classical family $SL(2,\mathbb{A})\cong\mathrm{Spin}(n+1,1)$ for the normed
division algebras $\mathbb{A}=\mathbb{R},\mathbb{C},\mathbb{H},\mathbb{O}$ of
dimension $n=1,2,4,8$:
\begin{equation}
\label{eq:app-lorentz:family}
\begin{aligned}
SL(2,\mathbb{R}) &\cong \mathrm{Spin}(2,1), &
SL(2,\mathbb{C}) &\cong \mathrm{Spin}(3,1), \\
SL(2,\mathbb{H}) &\cong \mathrm{Spin}(5,1), &
SL(2,\mathbb{O}) &\cong \mathrm{Spin}(9,1),
\end{aligned}
\end{equation}
with $\mathrm{Spin}(9,1)=\mathrm{Spin}(1,9)$ up to convention. \emph{[Theorem.]}
\end{theorem}

\begin{proof}[Proof (statement and status)]
The set of determinant- and Hermiticity-preserving actions on
$H_{2}(\mathbb{O})$ does not close under naive matrix multiplication, since
octonion non-associativity obstructs $g\,X\,g^{\dagger}$. Accordingly
$SL(2,\mathbb{O})$ is \emph{defined} as $\mathrm{Spin}(1,9)$ acting through
left/right octonion-multiplication generators chosen so that the determinant
\eqref{eq:app-lorentz:det} and the spinor bilinears $\psi^{\dagger}X\psi$ on
$\psi\in\mathbb{O}^{2}$ are covariant; the double cover
$\mathrm{Spin}(1,9)\to SO(1,9)$ is the statement that the spinor action is
two-valued on the vector action of \autoref{thm:app-lorentz:Mink}. These
isomorphisms are classical (Sudbery; Kugo--Townsend; Manogue--Schray; Baez).
\end{proof}

\begin{theorem}[The physical Lorentz four-plane; \texorpdfstring{$SL(2,\mathbb{C})\subset SL(2,\mathbb{O})$}{SL(2,C) in SL(2,O)}]
\label{thm:app-lorentz:SL2C}
The chain of composition-algebra inclusions
$\mathbb{R}\subset\mathbb{C}\subset\mathbb{H}\subset\mathbb{O}$ induces a chain of
spin inclusions; in particular
\begin{equation}
\label{eq:app-lorentz:c-in-o}
\mathrm{Spin}(1,3)=SL(2,\mathbb{C})\ \subset\ SL(2,\mathbb{O})=\mathrm{Spin}(1,9)
\end{equation}
via the sub-block $H_{2}(\mathbb{C})\subset H_{2}(\mathbb{O})$, where
$H_{2}(\mathbb{C})=\{[[a,z],[\bar{z},b]]:a,b\in\mathbb{R},\,z\in\mathbb{C}\}$ has
real dimension $4$, $\det=ab-|z|^{2}$ is the $(1,3)$ Minkowski form, and
$\det|_{H_{2}(\mathbb{C})}$ is $SL(2,\mathbb{C})$-covariant in the standard
four-dimensional way. The ambient $\mathbb{R}^{1,9}$ is the indefinite arena; the
\emph{physical} $(1,3)$ is a soldered four-plane inside it. \emph{[Theorem.]}
\end{theorem}

\begin{proof}
Restrict the change of variables of \autoref{thm:app-lorentz:Mink} to
$z\in\mathbb{C}\subset\mathbb{O}$: $\det=ab-|z|^{2}=(X^{0})^{2}-(X^{3})^{2}-(X^{1})^{2}-(X^{2})^{2}$
is the $(1,3)$ form. The covariance of the $SL(2,\mathbb{C})$ action on
$H_{2}(\mathbb{C})$ is the textbook 4D statement $X\mapsto gXg^{\dagger}$,
$\det(gXg^{\dagger})=|\det g|^{2}\det X=\det X$ for $g\in SL(2,\mathbb{C})$.
\end{proof}

\subsection{The cubic norm \texorpdfstring{$N$}{N} is genuinely indefinite}
\label{subsec:app-lorentz:cubic}

The Albert algebra $J_{3}(\mathbb{O})$ of $3\times3$ octonion-Hermitian matrices,
$\dim_{\mathbb{R}}=3+3\cdot8=27$, carries the quadratic trace form
$Q(X)=\mathrm{Tr}(X^{2})$ and the cubic norm $N(X)=\det X$
(see \autoref{thm:Invs} and \autoref{thm:OrbitClass}). On a diagonal Jordan frame
$X=\mathrm{diag}(\xi_{1},\xi_{2},\xi_{3})$ one has $Q=\sum_{i}\xi_{i}^{2}$ and
$N=\xi_{1}\xi_{2}\xi_{3}$.

\begin{theorem}[Separation of invariance groups]
\label{thm:app-lorentz:groups}
The automorphism group $F_{4}=\mathrm{Aut}(J_{3}(\mathbb{O}))$ (compact, $\dim 52$)
preserves the Jordan product, hence preserves \emph{both} $Q=\mathrm{Tr}(X^{2})$ and
$N$; since $Q$ is positive-definite, $F_{4}\subset SO(27)$. The reduced structure
group $E_{6(-26)}$ ($\dim 78$) preserves the cubic norm $N$ \emph{only up to a real
character} $\lambda(g)$, $N(gX)=\lambda(g)N(X)$, and does \emph{not} preserve $Q$.
\emph{[Theorem.]}
\end{theorem}

\begin{proof}
Automorphisms fix $\mathbb{1}$, commute with $\circ$, and hence fix every
coefficient of the characteristic polynomial, in particular $\mathrm{Tr}(X^{2})$ and
$\det=N$; Euclidean signature of $Q$ then forces $F_{4}\subset SO(27)$. The reduced
structure group is by definition the linear group fixing $N$ up to scalar. If it also
preserved $Q$ it would lie in $SO(27)$ and be compact, contradicting
$\dim E_{6(-26)}=78>52=\dim F_{4}$ together with the non-compactness of the real form
$E_{6(-26)}$ (its maximal compact is exactly $F_{4}$, so the $26$ non-compact
directions move $Q$). Hence $Q$ is not $E_{6}$-invariant; only $N$ is.
\end{proof}

\begin{remark}[Orbits are classified by rank and \texorpdfstring{$N$}{N}, not by \texorpdfstring{$(Q,N)$}{(Q,N)}]
\label{rmk:app-lorentz:orbits}
Real $E_{6}$-orbits on the $\mathbf{27}$ are classified by
$\mathrm{rank}\in\{0,1,2,3\}$ and the value of $N$, via the $E_{6}$-covariant
sharp map $X^{\#}$ ($\mathrm{rank}\,X\le1\Leftrightarrow X^{\#}=0$;
$\mathrm{rank}\,X\le2\Leftrightarrow N(X)=0$). It is \emph{not} the pair $(Q,N)$
that separates orbits, since $Q$ is $F_{4}$-invariant only. This corrects any
statement to the contrary (cf. \autoref{thm:OrbitClass}).
\end{remark}

\begin{theorem}[Indefiniteness, and the Minkowski block inside \texorpdfstring{$N$}{N}]
\label{thm:app-lorentz:indef}
The cubic norm $N$ takes strictly positive, strictly negative, and zero values on
$J_{3}(\mathbb{O})\cong\mathbb{R}^{27}$ (it is an odd-degree form,
$N(-X)=-N(X)$), so it is genuinely indefinite. Moreover, restricting $N$ to the
block obtained by fixing the third pivot idempotent, $\xi_{3}=1$, $z_{1}=z_{2}=0$,
and freezing $Y=\begin{psmallmatrix}\xi_{1}&z_{3}\\ \bar z_{3}&\xi_{2}\end{psmallmatrix}\in H_{2}(\mathbb{O})$,
\begin{equation}
\label{eq:app-lorentz:block}
N(X)\big|_{\mathrm{block}}=\xi_{1}\xi_{2}-n(z_{3})=\det Y,
\end{equation}
so the indefinite Minkowski form $-\det=-\eta$ of \autoref{thm:app-lorentz:Mink} is
the restriction of the cubic norm $N$ to an $H_{2}(\mathbb{O})\cong\mathbb{R}^{1,9}$
sub-block of $J_{3}(\mathbb{O})$. \emph{[Theorem.]}
\end{theorem}

\begin{proof}
$N(\mathrm{diag}(1,1,1))=1>0$, $N(\mathrm{diag}(-1,1,1))=-1<0$,
$N(\mathrm{diag}(0,1,1))=0$, and oddness is immediate from the cubic homogeneity;
\eqref{eq:app-lorentz:block} is the direct evaluation of $N$ on the stated block.
\end{proof}

\begin{theorem}[Consistency: definite \texorpdfstring{$K|_{\mathfrak{m}}$}{K|m} and indefinite \texorpdfstring{$N$}{N} live on different spaces]
\label{thm:app-lorentz:consistency}
The positive-definite coset/target metric $K|_{\mathfrak{m}}$ and the indefinite
spacetime form are mutually consistent because they are forms on \emph{different}
representation spaces. The coset metric is positive-definite on
$\mathfrak{m}=\mathbf{16}_{-3}\oplus\overline{\mathbf{16}}_{+3}$ (the adjoint
complement, the $32$-dimensional internal tangent of EIII), and is unique by Schur
($\mathfrak{m}$ is real-irreducible of complex type, commutant $\mathbb{C}$; see
\autoref{thm:uniqueMetric}). The cubic norm $N$ is indefinite and lives on the
$\mathbf{27}$, a \emph{different} module. Therefore ``the coset metric is Euclidean''
and ``the cubic norm is indefinite'' are statements about distinct bilinear data on
distinct modules; there is no tension. The Lorentzian signature used for spacetime is
imported from $N|_{H_{2}(\mathbb{O})}$ via the soldering map P3, and is \emph{never}
asked of $K|_{\mathfrak{m}}$. \emph{[Theorem.]}
\end{theorem}

\subsection{The soldering postulate (P3)}
\label{subsec:app-lorentz:p3}

We now state the soldering precisely. Throughout, P3 is a \emph{postulate}: added
structure, not a theorem of the coset geometry. The theorem-level inputs are: a
vacuum manifold $M=\mathrm{EIII}=E_{6}/(\mathrm{Spin}(10)\times U(1))$, $\dim_{\mathbb{R}}M=32$,
with reductive split $\mathfrak{e}_{6}=\mathfrak{h}\oplus\mathfrak{m}$,
$\mathfrak{h}=\mathfrak{so}(10)\oplus\mathfrak{u}(1)$,
$\mathfrak{m}=\mathbf{16}_{-3}\oplus\overline{\mathbf{16}}_{+3}$, and Maurer--Cartan
form $\theta=\theta_{\mathfrak{h}}+\theta_{\mathfrak{m}}$
(\autoref{cor:VacuumManifold}, \autoref{eq:MC}); together with the Minkowski-block
facts of \autoref{thm:app-lorentz:Mink}--\autoref{thm:app-lorentz:SL2C} and
\autoref{thm:app-lorentz:indef}. The $32$ clock fields are the coordinates along
$\mathfrak{m}$: $28$ are physical massive scalars (single common radiative mass,
Schur on the irreducible $\mathfrak{m}$) and $4$ are gauge/soldering modes eaten by
world-volume diffeomorphism plus local Lorentz (see \autoref{thm:MassSpectrum}).

\begin{postulate}[Soldering / frame identification, P3a]
\label{post:app-lorentz:P3a}
There exists a soldering map $\sigma$ identifying four distinguished clock-field
frame directions, written $e^{a}{}_{\mu}$ ($a=0,1,2,3$; $\mu$ a world index), with a
tangent four-plane $\Pi_{x}\subset T_{x}\mathcal{M}^{4}$ of an emergent four-manifold
$\mathcal{M}^{4}$. The vierbein is the pull-back of the Maurer--Cartan component
along these directions,
\begin{equation}
\label{eq:app-lorentz:vierbein}
e^{a}{}_{\mu}(x)=\big\langle\,\tau^{a},\,\theta_{\mathfrak{m}}(\partial_{\mu})\,\big\rangle,
\end{equation}
where $\{\tau^{a}\}_{a=0}^{3}$ is the image of the $H_{2}(\mathbb{C})\cong\mathbb{R}^{1,3}$
basis under the embedding of \autoref{post:app-lorentz:P3d}. The vierbein
\emph{emerges dynamically}, as the pull-back of $\theta_{\mathfrak{m}}$. \emph{[Postulate P3.]}
\end{postulate}

\begin{postulate}[Induced Minkowski form, P3b]
\label{post:app-lorentz:P3b}
The induced quadratic form on $\Pi_{x}$ is the Minkowski form inherited from the
cubic norm via the $H_{2}(\mathbb{C})$ block,
\begin{equation}
\label{eq:app-lorentz:metric}
g_{\mu\nu}(x)=e^{a}{}_{\mu}(x)\,e^{b}{}_{\nu}(x)\,\eta_{ab},
\qquad \eta_{ab}=\mathrm{diag}(-1,+1,+1,+1),
\end{equation}
with $\eta$ the restriction of $-\det=-N$ to the $H_{2}(\mathbb{C})$ sub-block
(\autoref{thm:app-lorentz:SL2C}, \autoref{thm:app-lorentz:indef}). \emph{[Postulate P3.]}
\end{postulate}

\begin{postulate}[Dimension and signature are added, P3c]
\label{post:app-lorentz:P3c}
The numbers $d=4$ and signature $(1,3)=(-,+,+,+)$ are \emph{part of P3}: added
structure. They are \emph{not} consequences of the Killing form, the coset metric, or
any rank lemma. \emph{[Postulate P3.]}
\end{postulate}

\begin{postulate}[Frame reduction by a vierbein condensate, P3d]
\label{post:app-lorentz:P3d}
A vierbein condensate (a VEV of the soldering bilinear) selects the
$H_{2}(\mathbb{C})\cong\mathbb{R}^{1,3}$ sub-block of
$H_{2}(\mathbb{O})\cong\mathbb{R}^{1,9}$ and breaks the ambient frame symmetry
\begin{equation}
\label{eq:app-lorentz:reduction}
\mathrm{Spin}(1,9)=SL(2,\mathbb{O})\ \longrightarrow\
SL(2,\mathbb{C})\times \mathrm{Spin}(6)=\mathrm{Spin}(1,3)\times\mathrm{Spin}(6),
\end{equation}
the orthogonal complement of $\mathbb{R}^{1,3}$ in $\mathbb{R}^{1,9}$ being the
six-dimensional spacelike $\mathbb{R}^{6}$ acted on by $\mathrm{Spin}(6)\cong SU(4)$
(internal). The residual $\mathrm{Spin}(1,3)=SL(2,\mathbb{C})$ is the local Lorentz
group of the soldered four-plane. \emph{[Postulate P3.]}
\end{postulate}

\begin{theorem}[Internal consistency of P3]
\label{thm:app-lorentz:P3consistency}
Given \autoref{post:app-lorentz:P3a}--\autoref{post:app-lorentz:P3d}, the tensor
$g_{\mu\nu}$ of \eqref{eq:app-lorentz:metric} is a non-degenerate symmetric rank-$2$
tensor of Lorentzian signature $(1,3)$ on $\mathcal{M}^{4}$, invariant under the
residual local Lorentz group $SO(1,3)$; the four frame modes $e^{a}{}_{\mu}$ carry
exactly the gauge degrees of freedom eaten by world-volume diffeomorphisms together
with local Lorentz, consistent with the count $32=28\ \text{physical}+4\ \text{gauge/soldering}$.
\emph{[Theorem (conditional on P3).]}
\end{theorem}

\begin{proof}
Non-degeneracy and signature are inherited from
$\eta_{ab}=\mathrm{diag}(-1,+1,+1,+1)$ (\autoref{thm:app-lorentz:SL2C}) whenever
$\det(e^{a}{}_{\mu})\ne0$, which holds for a generic condensate;
$SO(1,3)$-invariance is the defining relation
$\eta_{ab}=\Lambda^{c}{}_{a}\Lambda^{d}{}_{b}\eta_{cd}$ for
$\Lambda\in SO(1,3)$. The gauge count is the Ward identity
$\nabla_{\mu}\big(\delta\Gamma/\delta e^{a}{}_{\mu}\big)=0$, so the four frame
directions are protected by diffeomorphism plus local Lorentz, not by a Hessian zero.
\end{proof}

\begin{remark}[What P3 does not claim]
\label{rmk:app-lorentz:notclaim}
P3 does not claim that $K|_{\mathfrak{m}}$ is Lorentzian: it stays positive-definite,
the kinetic/target-space metric of the $\sigma$-model. P3 does not use a $(4,28)$
split of $\mathfrak{m}$ with Lorentzian signature: the $4$ are gauge soldering modes,
the $28$ are Euclidean target coordinates. P3 does not build the four-plane from
$\nabla I_{2},\nabla I_{3}$ singlets: $\mathfrak{m}$ contains no $SO(10)$ singlet, and
$\nabla Q|_{X_{0}}=2X_{0}\notin\mathfrak{m}$. The four-plane is the soldered image of
$H_{2}(\mathbb{C})\subset\mathbf{27}$, not a sub-module of $\mathfrak{m}$.
\end{remark}

\begin{remark}[Anti-theorem: there is no Killing-form derivation of \texorpdfstring{$(1,3)$}{(1,3)}]
\label{rmk:app-lorentz:anti}
There is \emph{no} derivation of $d=4$ or of the signature $(1,3)$ from the Killing
form of $E_{6}$, from the coset metric $K|_{\mathfrak{m}}$, or from any rank lemma.
The Killing form of compact $E_{6}$ is negative-definite; its restriction to
$\mathfrak{m}$ is (sign-flipped to) positive-definite of signature $(32,0)$ --- zero
Lorentzian content. Any claim of a $(4,28)$ or $(1,3)$ Killing/coset signature is
false. Lorentzian signature enters only through P3, rooted in the indefinite cubic
norm $N$ / $H_{2}(\mathbb{O})\cong\mathbb{R}^{1,9}$.
\end{remark}

\subsection{Emergent time: Page--Wootters clock and Osterwalder--Schrader reconstruction}
\label{subsec:app-lorentz:time}

Lorentzian time is recovered by a relational-clock plus reconstruction-theorem
mechanism; no ``multiply $e^{0}$ by $i$'' Wick trick is used. The soldered metric
$\eta^{(1,3)}$ is already Lorentzian \emph{before} any continuation; the role of
analytic continuation is confined to the correlation functions in the clock-time
argument, justified by reflection positivity.

\begin{postulate}[Timelike clock selection, P3e]
\label{post:app-lorentz:clock}
Among the $32$ clock fields, the soldering condensate
(\autoref{post:app-lorentz:P3d}) selects a distinguished timelike direction, whose
clock field we call $\phi^{0}$ --- concretely the trace/$X^{0}$ direction of the
soldered $H_{2}(\mathbb{C})\subset J_{3}(\mathbb{O})$ block (the unique
$\eta$-timelike soldering leg $e^{0}{}_{\mu}$). The selection of \emph{one} timelike
direction is precisely the $(1,3)$ (rather than $(2,8)$, etc.) content of P3.
\emph{[Postulate P3.]}
\end{postulate}

\begin{theorem}[Page--Wootters relational evolution]
\label{thm:app-lorentz:PW}
Promote $\phi^{0}$ to a Page--Wootters clock: enlarge the kinematical Hilbert space
to $\mathcal{H}=\mathcal{H}_{\mathrm{clock}}\otimes\mathcal{H}_{\mathrm{sys}}$ and
impose the induced-gravity Hamiltonian constraint
\begin{equation}
\label{eq:app-lorentz:constraint}
\hat{H}_{\mathrm{phys}}\,|\Psi\rangle\!\rangle=0,
\qquad
\hat{H}_{\mathrm{phys}}=\hat{H}_{\mathrm{clock}}+\hat{H}_{\mathrm{sys}}
\ (+\ \text{interaction}),
\end{equation}
on physical states. Conditioning on clock readings $\phi^{0}=\tau$,
$|\psi(\tau)\rangle_{\mathrm{sys}}=\langle\phi^{0}=\tau|\Psi\rangle\!\rangle$, yields
the Schr\"odinger equation in relational time,
\begin{equation}
\label{eq:app-lorentz:schrodinger}
i\,\partial_{\tau}|\psi(\tau)\rangle_{\mathrm{sys}}
 =\hat{H}_{\mathrm{sys}}|\psi(\tau)\rangle_{\mathrm{sys}},
\end{equation}
with the unitary group $U(\tau)=e^{-i\hat{H}_{\mathrm{sys}}\tau}$ recovered on
physical states in the semiclassical regime where the soldered $g_{\mu\nu}$ admits a
timelike Killing vector and $\hat{H}_{\mathrm{sys}}=\int T_{00}$ along it.
\emph{[Theorem (Page--Wootters), conditional on the constraint and a good clock; global hyperbolicity assumed.]}
\end{theorem}

\begin{proof}
Under \eqref{eq:app-lorentz:constraint} with a good clock
($\hat{H}_{\mathrm{clock}}=\hat{p}_{\phi^{0}}$ canonically conjugate to
$\hat{\phi}^{0}$), the conditioned state obeys \eqref{eq:app-lorentz:schrodinger};
this is the Page--Wootters theorem. The identification $t\leftrightarrow X^{0}$ is
$\hat{H}_{\mathrm{sys}}=\int T_{00}$ along the timelike Killing direction.
\end{proof}

\begin{construction}[Osterwalder--Schrader reconstruction along \texorpdfstring{$\phi^{0}$}{phi0}]
\label{con:app-lorentz:OS}
The Euclidean $\sigma$-model built on the positive-definite $K|_{\mathfrak{m}}$
defines a Euclidean functional measure with Schwinger functions
$S_{n}(x_{1},\dots,x_{n})$. Take $\tau=\phi^{0}$ as Osterwalder--Schrader time and
assume the axioms: (OS0) Euclidean invariance of the spatial slice; (OS1) reflection
positivity with respect to $\theta:\tau\mapsto-\tau$ about a $\phi^{0}$-slice,
\begin{equation}
\label{eq:app-lorentz:rp}
\sum_{i,j}\bar{c}_{i}\,c_{j}\,S\!\left(\theta F_{i}\cdot F_{j}\right)\ \ge\ 0,
\qquad \mathrm{supp}\,F_{i}\subset\{\tau>0\};
\end{equation}
(OS2) symmetry and clustering.
\end{construction}

\begin{theorem}[OS reconstruction \texorpdfstring{$\Rightarrow$}{=>} Lorentzian Wightman theory]
\label{thm:app-lorentz:OS}
Given (OS0)--(OS2), the Osterwalder--Schrader reconstruction theorem produces a
Hilbert space $\mathcal{H}$, a positive self-adjoint Hamiltonian $\hat{H}\ge0$
generating $\tau$-translations as the contraction semigroup $e^{-\hat{H}\tau}$, and,
by analytic continuation of the Schwinger functions in the time argument
$\tau\to it$, Lorentzian Wightman functions on $(\mathcal{M}^{4},\eta^{(1,3)})$ with
unitary evolution $e^{-i\hat{H}t}$ covariant under the reconstructed $SO(1,3)$. The
OS theorem itself is rigorous given its hypotheses \emph{[Theorem]}; that the UCFT
$\sigma$-model \emph{satisfies} reflection positivity \eqref{eq:app-lorentz:rp} along
$\phi^{0}$ is \emph{[Indicative]} --- the precise property to be checked, and the
rigorous replacement for the Wick trick.
\end{theorem}

\begin{proof}[Proof (theorem part)]
Reflection positivity \eqref{eq:app-lorentz:rp} induces a physical inner product on
the quotient $\{F:\tau>0\}/(\text{null})$; the generator of the resulting contraction
semigroup $e^{-\hat{H}\tau}$ is self-adjoint and positive. The Schwinger functions,
analytic in a tube, continue to Wightman distributions satisfying the Wightman axioms,
including the spectrum condition and Lorentz covariance under the reconstructed
$SO(1,3)$ (Osterwalder--Schrader, 1973, 1975).
\end{proof}

\begin{remark}[Why this is not the Wick trick]
\label{rmk:app-lorentz:notwick}
The continuation $\tau\to it$ acts on the correlation functions in the
\emph{clock-time argument}, justified by reflection positivity
\eqref{eq:app-lorentz:rp} and the resulting analyticity tube --- it is \emph{not} the
ad hoc replacement $e^{0}{}_{\mu}\mapsto i\,e^{0}{}_{\mu}$ of a frame field. The
metric $\eta^{(1,3)}$ is the soldered (P3) form, Lorentzian before any continuation;
OS reconstruction supplies the unitary dynamics $e^{-i\hat{H}t}$ and the
positive-energy spectrum, with $\phi^{0}$ (the Page--Wootters clock) serving as the
OS reflection time. The two mechanisms dovetail:
\autoref{thm:app-lorentz:PW} gives relational unitarity $U(\tau)$ on physical states,
and \autoref{thm:app-lorentz:OS} gives the positive-energy, Lorentz-covariant
Wightman theory.
\end{remark}

\subsection{Summary: theorem versus postulate}
\label{subsec:app-lorentz:summary}

The division of labour is as follows.
\begin{itemize}
\item \textbf{Theorems} (proved): $(H_{2}(\mathbb{O}),-\det)\cong\mathbb{R}^{1,9}$
  (\autoref{thm:app-lorentz:Mink}); $SL(2,\mathbb{O})\cong\mathrm{Spin}(1,9)$ and
  $SL(2,\mathbb{C})\cong\mathrm{Spin}(1,3)$
  (\autoref{thm:app-lorentz:SL2O}, \autoref{thm:app-lorentz:SL2C}); $F_{4}$ preserves
  $Q$, $E_{6}$ preserves only $N$ (\autoref{thm:app-lorentz:groups}); $N$ indefinite
  and the Minkowski block $\det Y\subset N$ (\autoref{thm:app-lorentz:indef});
  definite $K|_{\mathfrak{m}}$ and indefinite $N$ are consistent
  (\autoref{thm:app-lorentz:consistency}); internal consistency of $g_{\mu\nu}$ given
  P3 (\autoref{thm:app-lorentz:P3consistency}); the Page--Wootters and OS
  reconstruction theorems (\autoref{thm:app-lorentz:PW}, \autoref{thm:app-lorentz:OS}).
\item \textbf{Postulate} (P3, added structure, cf. P1--P6): the soldering map, the
  choice of a four-plane, $d=4$, signature $(1,3)$, the vierbein condensate and frame
  reduction, and the timelike clock selection
  (\autoref{post:app-lorentz:P3a}--\autoref{post:app-lorentz:P3d},
  \autoref{post:app-lorentz:clock}).
\item \textbf{Indicative}: that the UCFT $\sigma$-model satisfies reflection
  positivity along $\phi^{0}$ (\autoref{thm:app-lorentz:OS}).
\end{itemize}
In short, the Lorentzian arena is a \emph{conditional embedding}: the algebra
$H_{2}(\mathbb{O})\cong\mathbb{R}^{1,9}$ and $SL(2,\mathbb{O})\cong\mathrm{Spin}(1,9)$
are theorems, while the selection of the physical $(1,3)$ four-plane and the recovery
of unitary Lorentzian dynamics rest on the soldering postulate P3 plus an indicative
reflection-positivity assumption.
